\documentclass[11pt]{article}
\usepackage[margin=1in]{geometry}
\usepackage{amssymb,stmaryrd}
\usepackage{enumitem}
\usepackage{xcolor}
\usepackage{natbib,bibentry}
\usepackage[hidelinks]{hyperref}
\bibliographystyle{plainnat}

% Status of a task, drawn like the [ ] [/] [x] [-] of our plain-text task lists.
\newcommand{\todo}{$\square$}
\newcommand{\doing}{$\boxslash$}
\newcommand{\done}{$\boxtimes$}
\newcommand{\dropped}{$\boxminus$}
% A field not written yet.
\newcommand{\blank}{\textcolor{gray}{---}}
% A paper in the reading list: its full reference from Roadmap.bib, then why we read it.
\newcommand{\paper}[2]{\bibentry{#1}\par\textit{#2}}
% Marks a paper to read first.
\newcommand{\first}{$\bigstar$\ }

% An aim's fields: the name on the left, the text beside it.
\newenvironment{fields}
  {\begin{description}[style=sameline, leftmargin=7em, labelwidth=6.5em, font=\normalfont\bfseries]}
  {\end{description}}
% A list of tasks; each \item takes its status as the label, e.g. \item[\todo].
\newenvironment{tasks}
  {\begin{itemize}[leftmargin=1.5em, itemsep=2pt]}
  {\end{itemize}}

\setlength{\parindent}{0pt}
\setlength{\parskip}{4pt}

\begin{document}
% Reads the references for \bibentry without printing a bibliography; \nocite takes every entry.
\nobibliography{Roadmap}
\nocite{*}

\begin{center}
  {\Large\bfseries Project Roadmap}\\[4pt]
  Started \blank \qquad
  \todo\ to do \quad \doing\ doing \quad \done\ done \quad \dropped\ dropped
\end{center}

\section*{Central question}
\blank

\section*{Order}
\begin{itemize}[leftmargin=1.5em, itemsep=2pt]
  \item \blank
  \item Within the aims, start with the cheap step that is most likely to fail.
  \item For each experiment, fix what it compares, its metric and its output before writing code.
\end{itemize}

% One section per aim; copy it for Aim B, C, ...
\section*{Aim A. \blank}
\begin{fields}
  \item[Hypothesis] \blank
  \item[Theory] \todo\ \blank
  \item[Experiment] \todo\ \blank
  \item[Outcomes] \blank
  \item[Risks] \blank
  \item[Done when] \blank
  \item[Paper] undecided
\end{fields}

\section*{Templates}
\begin{tasks}
  \item[\todo] Poster.
  \item[\todo] The rest of the template (Research Template in Projects.md): arXiv, project page, code,
    demo video, model card, slides, video, the research process and how we work with AI.
\end{tasks}
\begin{fields}
  \item[Done when] \blank
\end{fields}

\section*{Code}
Make the code easier to read, and keep the experiments apart.
\begin{tasks}
  \item[\todo] One command reproduces each figure in \path{figs/}, with its data source written down.
\end{tasks}
\begin{fields}
  \item[Done when] Every figure in \path{figs/} comes out of a clean clone.
\end{fields}

\section*{Reading}
\begin{fields}
  \item[Topics] \blank
  \item[Done when] A list with one sentence per paper on its relation to this project.
\end{fields}

\section*{Open decisions}
\begin{itemize}[leftmargin=1.5em, itemsep=2pt]
  \item \blank
\end{itemize}

% One subsection per category. Add a paper: its BibTeX (bib-lookup skill) to Roadmap.bib,
% and \item[\todo] \paper{key}{Why we read it.} under its category.
\section*{Reading list}
\subsection*{\blank}
\begin{tasks}
  \item[\todo] \blank
\end{tasks}

\end{document}
